\documentclass[11pt,a4paper]{article}
\usepackage[utf8]{inputenc}
\usepackage{amsmath,amssymb,amsthm}
\usepackage[margin=1in]{geometry}

\newtheorem{theorem}{Theorem}
\newtheorem{lemma}[theorem]{Lemma}

\title{Problem 130: Composites with Prime Repunit Property}
\author{Project Euler Solution}
\date{}

\begin{document}
\maketitle

\section{Problem Statement}

Define $R(k)$ as the repunit of length $k$ and $A(n)$ as the least $k$ such that $n \mid R(k)$. For all primes $p > 5$, $A(p) \mid (p-1)$. Find the sum of the first 25 composite $n$ with $\gcd(n,10) = 1$ and $A(n) \mid (n-1)$.

\section{Mathematical Foundation}

\begin{theorem}[Prime repunit property]
For any prime $p > 5$, $A(p) \mid (p-1)$.
\end{theorem}

\begin{proof}
By Fermat's Little Theorem, $10^{p-1} \equiv 1 \pmod{p}$. Since $\gcd(9, p) = 1$:
\[
R(p-1) = \frac{10^{p-1} - 1}{9} \equiv 0 \pmod{p}.
\]
Thus $p \mid R(p-1)$. Since $A(p) = \mathrm{ord}_p(10)$ and $\mathrm{ord}_p(10) \mid (p-1)$, we get $A(p) \mid (p-1)$.
\end{proof}

\begin{theorem}[Pseudoprime connection]
If $n$ is composite with $\gcd(n, 10) = 1$, $\gcd(n, 9) = 1$, and $10^{n-1} \equiv 1 \pmod{n}$, then $A(n) \mid (n-1)$.
\end{theorem}

\begin{proof}
Since $\gcd(n,9) = 1$, $A(n) = \mathrm{ord}_n(10)$. If $10^{n-1} \equiv 1 \pmod{n}$, then $\mathrm{ord}_n(10) \mid (n-1)$.
\end{proof}

\begin{lemma}[Divisibility]
If $n \mid R(k)$, then $A(n) \mid k$.
\end{lemma}

\begin{proof}
The sequence $R(k) \bmod n$ is periodic with period $A(n)$, and $R(A(n)) \equiv 0$. Thus $R(k) \equiv 0 \pmod{n}$ iff $A(n) \mid k$.
\end{proof}

\begin{lemma}[Iterative computation]
$A(n)$ is computed via $r_1 = 1 \bmod n$, $r_k = (10r_{k-1} + 1) \bmod n$. The first $k$ with $r_k = 0$ is $A(n)$.
\end{lemma}

\begin{proof}
Follows from $R(k) = 10R(k-1) + 1$.
\end{proof}

\section{Algorithm}

\begin{enumerate}
\item For $n = 2, 3, 4, \ldots$:
\begin{enumerate}
\item Skip if $\gcd(n, 10) \ne 1$ or $n$ is prime.
\item Compute $A(n)$ via the recurrence.
\item If $(n-1) \bmod A(n) = 0$, add $n$ to the list.
\end{enumerate}
\item Stop after 25 composites. Return their sum.
\end{enumerate}

\section{Complexity Analysis}

\begin{itemize}
\item \textbf{Time:} $O(A(n)) \le O(n)$ per candidate for computing $A(n)$. $O(\sqrt{n})$ for primality test. The 25 composites are found among small values, so total work is modest.
\item \textbf{Space:} $O(1)$.
\end{itemize}

\section{Answer}

\[
\boxed{149253}
\]

\end{document}
